\documentclass[10pt, letterpaper]{article}

% Packages:
\usepackage[
    ignoreheadfoot,
    top=1.5cm,
    bottom=1.5cm,
    left=1.5cm,
    right=1.5cm,
    footskip=0.7cm,
]{geometry}
\usepackage{titlesec} % for customizing section titles
\usepackage{tabularx} % for making tables with fixed width columns
\usepackage{array} % tabularx requires this
\usepackage[dvipsnames]{xcolor} % for coloring text
\definecolor{primaryColor}{RGB}{0, 0, 0} % define primary color
\usepackage{enumitem} % for customizing lists
\usepackage{fontawesome5} % for using icons
\usepackage{amsmath} % for math
\usepackage[
    pdftitle={CV},
    pdfauthor={Your Name},
    colorlinks=true,
    urlcolor=primaryColor
]{hyperref}
\usepackage[pscoord]{eso-pic} % for floating text on the page
\usepackage{calc} % for calculating lengths
\usepackage{bookmark} % for bookmarks
\usepackage{lastpage} % for getting the total number of pages
\usepackage{changepage} % for one column entries (adjustwidth environment)
\usepackage{paracol} % for two and three column entries
\usepackage{ifthen} % for conditional statements
\usepackage{needspace} % for avoiding page brake right after the section title
\usepackage{iftex} % check if engine is pdflatex, xetex or luatex

% Ensure that generate pdf is machine readable/ATS parsable:
\ifPDFTeX
    \input{glyphtounicode}
    \pdfgentounicode=1
    \usepackage[T1]{fontenc}
    \usepackage[utf8]{inputenc}
    \usepackage{lmodern}
\fi

\usepackage{charter}

% Some settings:
\raggedright
\AtBeginEnvironment{adjustwidth}{\partopsep0pt} % remove space before adjustwidth environment
\pagestyle{empty} % no header or footer
\setcounter{secnumdepth}{0} % no section numbering
\setlength{\parindent}{0pt} % no indentation
\setlength{\topskip}{0pt} % no top skip
\setlength{\columnsep}{0.15cm} % set column seperation
\pagenumbering{gobble} % no page numbering

\titleformat{\section}{\needspace{4\baselineskip}\bfseries\large}{}{0pt}{}[\vspace{1pt}\titlerule]

\titlespacing{\section}{
    % left space:
    -1pt
}{
    % top space:
    0.3 cm
}{
    % bottom space:
    0.2 cm
} % section title spacing

\renewcommand\labelitemi{$\vcenter{\hbox{\small$\bullet$}}$} % custom bullet points
\newenvironment{highlights}{
    \begin{itemize}[
        topsep=0.10 cm,
        parsep=0.10 cm,
        partopsep=0pt,
        itemsep=0pt,
        leftmargin=0 cm + 10pt
    ]
}{
    \end{itemize}
} % new environment for highlights


\newenvironment{highlightsforbulletentries}{
    \begin{itemize}[
        topsep=0.10 cm,
        parsep=0.10 cm,
        partopsep=0pt,
        itemsep=0pt,
        leftmargin=10pt
    ]
}{
    \end{itemize}
} % new environment for highlights for bullet entries

\newenvironment{onecolentry}{
    \begin{adjustwidth}{
        0 cm + 0.00001 cm
    }{
        0 cm + 0.00001 cm
    }
}{
    \end{adjustwidth}
} % new environment for one column entries

\newenvironment{twocolentry}[2][]{
    \onecolentry
    \def\secondColumn{#2}
    \setcolumnwidth{\fill, 4.5 cm}
    \begin{paracol}{2}
}{
    \switchcolumn \raggedleft \secondColumn
    \end{paracol}
    \endonecolentry
} % new environment for two column entries

\newenvironment{threecolentry}[3][]{
    \onecolentry
    \def\thirdColumn{#3}
    \setcolumnwidth{, \fill, 4.5 cm}
    \begin{paracol}{3}
    {\raggedright #2} \switchcolumn
}{
    \switchcolumn \raggedleft \thirdColumn
    \end{paracol}
    \endonecolentry
} % new environment for three column entries

\newenvironment{header}{
    \setlength{\topsep}{0pt}\par\kern\topsep\centering\linespread{1.5}
}{
    \par\kern\topsep
} % new environment for the header

% save the original href command in a new command:
\let\hrefWithoutArrow\href

\begin{document}
    \newcommand{\AND}{\unskip
        \cleaders\copy\ANDbox\hskip\wd\ANDbox
        \ignorespaces
    }
    \newsavebox\ANDbox
    \sbox\ANDbox{$|$}

    \begin{header}
        \fontsize{20 pt}{20 pt}\selectfont Your Name

        \vspace{3 pt}

        \normalsize
        \mbox{Your City, Country}%
        \kern 5.0 pt%
        \AND%
        \kern 5.0 pt%
        \mbox{\hrefWithoutArrow{mailto:you@example.com}{you@example.com}}%
        \kern 5.0 pt%
        \AND%
        \kern 5.0 pt%
        \mbox{\hrefWithoutArrow{tel:+1-234-567-8900}{+1-234-567-8900}}%
        \kern 5.0 pt%
        \AND%
        \kern 5.0 pt%
        \mbox{\hrefWithoutArrow{https://linkedin.com/in/your-handle}{linkedin.com/in/your-handle}}%
        \kern 5.0 pt%
        \AND%
        \kern 5.0 pt%
        \mbox{\hrefWithoutArrow{https://github.com/octocat}{github.com/octocat}}%
    \end{header}

    \vspace{5 pt - 0.3 cm}

\hrulefill

\begin{onecolentry}
    \textbf{Summary-} A short professional summary. Describe who you are, your focus areas, and the kind of impact you make. Replace this with your own.
\end{onecolentry}


\section{Skills}

\begin{onecolentry}
    \textbf{Languages:} Python, JavaScript, TypeScript, Java, HTML, CSS
\end{onecolentry}

\vspace{0.1 cm}

\begin{onecolentry}
    \textbf{Frameworks \& Libraries:} Django, FastAPI, ReactJS, NodeJS
\end{onecolentry}

\vspace{0.1 cm}

\begin{onecolentry}
    \textbf{Tools \& Technologies:} Linux, Git, PostgreSQL, MongoDB, AWS, Docker
\end{onecolentry}


\section{Work Experience}


\begin{twocolentry}{
    Jan. 2024 – Present
}
    \textbf{Software Engineer}, Acme Corp, City
\end{twocolentry}
\begin{onecolentry}
    \begin{highlights}
        \item Describe a key responsibility or achievement, including the tech you used and the outcome.
        \item Quantify impact where you can (e.g. "reduced load time by 40\%").
    \end{highlights}
\end{onecolentry}

\vspace{0.2 cm}

\begin{twocolentry}{
    Jun. 2022 – Dec. 2023
}
    \textbf{Full-Stack Developer}, Startup Inc, City
\end{twocolentry}
\begin{onecolentry}
    \begin{highlights}
        \item Built and shipped features across the stack.
        \item Designed REST APIs and integrated third-party services.
    \end{highlights}
\end{onecolentry}


\section{Education}

\begin{onecolentry}
    \textbf{Your University}, Bachelor of Science in Computer Science
\end{onecolentry}

\vspace{0.1 cm}

\begin{onecolentry}
    \textbf{Your High School}, A Levels
\end{onecolentry}


\section{Personal Projects}

\begin{onecolentry}
     \textbf{Project One} – \textit{\href{https://github.com/octocat}{Repo}}
\end{onecolentry}
\begin{onecolentry}
    \begin{highlights}
        \item One or two lines on what the project does and the stack it uses.
    \end{highlights}
\end{onecolentry}

\vspace{0.2 cm}

\begin{onecolentry}
    \textbf{\href{https://github.com/octocat}{Project Two (github.com/you/project-two)}}
\end{onecolentry}
\begin{onecolentry}
    \begin{highlights}
        \item Describe the project, the problem it solves, and your role.
    \end{highlights}
\end{onecolentry}


\section{Volunteering and Related Experience}

\begin{onecolentry}
    \textbf{Volunteer Role}, Some Community / Club
\end{onecolentry}
\begin{highlights}
\item Describe what you did and the impact you had.
\end{highlights}


\end{document}
